\documentclass[12pt]{article}
\usepackage[margin=1in]{geometry}
\usepackage{amsmath,amssymb}
\usepackage{graphicx}
\usepackage{booktabs}
\usepackage{hyperref}
\usepackage{natbib}
\usepackage{lineno}
\usepackage{setspace}

\doublespacing

\title{An Allocentric Human Odometer for Perceiving Distances on the Ground Plane}

\author{Author One$^{1}$, Author Two$^{1}$, Author Three$^{1}$ \\
\small $^{1}$Department of Psychological and Brain Sciences, University of Example, City, Country}

\date{}

\begin{document}

\maketitle

\begin{abstract}
How the visual system maintains stable spatial representations during locomotion remains a fundamental question in perception science. Prior studies of visual space perception tested only stationary observers, leaving open whether spatial coding operates in an allocentric (world-centered) or egocentric (body-centered) reference frame. Here, across five experiments, we demonstrate that the intrinsic bias---an implicit curved surface representation used to localize objects in the dark---is anchored allocentrically at the observer's starting location rather than traveling with the observer during forward locomotion. We show that this allocentric anchoring decays over time, requires attentional resources during walking, depends on vestibular signals (passive wheelchair transport is sufficient while active proprioception is not necessary), is sensitive to the direction of travel (backward transport reverses the effect), and operates anisotropically along the horizontal direction only. These findings reveal a human path-integration mechanism, analogous to the horizontal-only odometer described in desert ants, that maintains allocentric spatial references through vestibular-driven distance monitoring during locomotion.
\end{abstract}

\section{Introduction}

A remarkable feature of visual perception is its apparent stability despite the continuous retinal image motion caused by self-movement through the environment. As we walk, our eyes, head, and body translate and rotate, producing a constantly changing retinal image, yet our perception of the spatial layout of objects remains largely veridical \citep{gibson1979, loomis2002}. Understanding the coordinate system that supports this perceptual stability is essential for a comprehensive account of spatial vision.

Two fundamental coordinate systems have been proposed for spatial representation: egocentric (body-centered) and allocentric (world-centered) \citep{klatzky1998, burgess2006}. In an egocentric frame, spatial positions are coded relative to the observer's current position, so that the representation must be continuously updated as the observer moves. In an allocentric frame, positions are coded relative to fixed landmarks or environmental features, providing a stable reference that does not require updating during locomotion. While both types of representation have been demonstrated in the context of spatial memory and navigation \citep{mou2004, waller2006}, their role in visual space perception per se has been difficult to disentangle because, for a stationary observer, both coordinate systems predict identical perceptual outcomes.

An important feature of visual space perception is the intrinsic bias, an implicit curved surface representation that the visual system uses to localize objects in the dark \citep{ooi2001, wu2004}. When observers judge the position of a briefly presented target in an otherwise dark environment, their responses are consistent with a representation of the ground plane that curves upward at far distances, as though the visual system assumes a concave surface in the absence of visible texture gradients or other depth cues. When visible ground texture is available, the intrinsic bias integrates with these cues to produce more accurate spatial estimates \citep{he2004}. The intrinsic bias has been characterized in stationary observers, but whether it remains fixed in world coordinates or moves with the observer during locomotion has not been tested.

The distinction between allocentric and egocentric coding of the intrinsic bias can be revealed by having observers walk in the dark before making spatial judgments. If the intrinsic bias is coded allocentrically, it should remain anchored at the observer's starting position (the ``home base'') after forward walking. Consequently, targets presented from the new, more distant position should appear nearer than they would from the home base, because the observer has moved forward relative to the fixed curved surface representation. Alternatively, if the intrinsic bias is coded egocentrically, it should move with the observer, and targets should appear at the same locations regardless of whether the observer has walked.

Here we report five experiments that test these predictions and characterize the path-integration mechanism underlying allocentric coding. Experiment~1 establishes that the intrinsic bias remains at the home base after forward walking, supporting allocentric coding, and shows that this effect decays over time. Experiment~2 demonstrates that dividing attention during walking disrupts path-integration, causing the intrinsic bias to travel with the observer. Experiment~3 shows that passive vestibular translation via wheelchair transport is sufficient for path-integration, while Experiment~4 confirms that backward translation shifts the bias in the opposite direction. Experiment~5 reveals that the path-integration mechanism operates anisotropically, gauging only horizontal distance and not vertical displacement.

\section{Methods}

\subsection{Participants}

Fourteen observers (7 male, 7 female; mean age $24.47 \pm 1.45$ years; mean eye height $1.60 \pm 0.02$~m) participated in Experiments 1--4, with 8 observers per experiment in a within-subject design (Table~\ref{tab:participants}). Nine additional observers (6 male, 3 female; mean age $22.78 \pm 1.22$ years; mean eye height $1.58 \pm 0.04$~m) participated in Experiment~5. All observers had visual acuity of at least 20/20 and stereoscopic resolution of 20 arc seconds or better. The study was approved by the institutional review board and all participants gave informed consent.

\begin{table}[ht]
\centering
\caption{Participant demographics and experimental design across all five experiments.}
\label{tab:participants}
\begin{tabular}{lcccccc}
\toprule
\textbf{Parameter} & \textbf{Exp 1} & \textbf{Exp 2} & \textbf{Exp 3} & \textbf{Exp 4} & \textbf{Exp 5} \\
\midrule
N & 8 & 8 & 8 & 8 & 9 \\
Design & Within-subj & Within-subj & Within-subj & Within-subj & Within-subj \\
Mean age (yr) & \multicolumn{4}{c}{$24.47 \pm 1.45$} & $22.78 \pm 1.22$ \\
Mean eye height (m) & \multicolumn{4}{c}{$1.60 \pm 0.02$} & $1.58 \pm 0.04$ \\
Sex (M/F) & \multicolumn{4}{c}{7/7 (total)} & 6/3 \\
\bottomrule
\end{tabular}
\end{table}

\subsection{Apparatus and Stimuli}

All experiments were conducted in a large dark room whose layout was unknown to observers. The target was a diffused green LED mounted inside a ping-pong ball (luminance 0.16 cd/m$^2$; angular size 0.22$^\circ$ at eye level), presented at 5 Hz flicker for 1 second. Music was played through speakers throughout the session to mask auditory cues that might provide spatial information. A guidance rope was mounted at 0.8~m above the floor to guide observers during blind walking.

The testing area was separated from a dimly lit waiting area (approximately 3~m$^2$) by a black curtain. Between trials, observers sat in the waiting area facing a wall, with their back to the testing area, preventing them from learning the room layout. A piece of plastic wrap on the guidance rope marked the starting position (home base) for each trial.

\subsection{Response Task}

Observers performed a blind walking-gesturing task. After viewing the target for 1 second, the experimenter removed the target and the observer walked blindly while sliding one hand along the guidance rope to the remembered target location, then gestured the perceived height of the target. The experimenter shook the rope as a signal when it was safe to walk. This response task provides a measure of the observer's perceived target location in both horizontal distance and vertical height.

\subsection{Experiment 1: Allocentric vs.\ Egocentric Coding}

Three conditions were tested (Table~\ref{tab:exp1conditions}). In the baseline-stationary condition, observers stood at the home base, waited 12 or 60 seconds, and then judged the target location. In the walking conditions, observers walked 1.5~m forward from the home base in the dark, waited either 12 or 60 seconds, and then judged the target location. The allocentric hypothesis predicts that walking conditions should produce nearer judged locations (because the intrinsic bias remains at the home base), while the egocentric hypothesis predicts no difference.

\begin{table}[ht]
\centering
\caption{Experimental conditions and key results for Experiment 1.}
\label{tab:exp1conditions}
\begin{tabular}{llll}
\toprule
\textbf{Condition} & \textbf{Delay} & \textbf{Manipulation} & \textbf{Result} \\
\midrule
Baseline-stationary & 12 s & Stand at home base & Reference \\
Walking (12 s) & 12 s & Walk 1.5 m forward & Judged locations nearer (significant) \\
Walking (60 s) & 60 s & Walk 1.5 m forward & Smaller shift (effect decays) \\
\bottomrule
\end{tabular}
\end{table}

\subsection{Experiment 2: Divided Attention During Walking}

Two conditions were tested: baseline-stationary and divided-attention-walking. In the divided-attention condition, observers walked 1.5~m forward while performing a demanding cognitive task (counting backward by 7s). If path-integration requires attentional resources, then dividing attention should prevent the allocentric anchoring of the intrinsic bias, and judged locations should be similar to baseline.

\subsection{Experiment 3: Passive Vestibular Forward Translation}

Three conditions were tested: baseline-stationary, active walking (1.5~m forward), and vestibular-forward (1.5~m passive wheelchair transport forward). Observers were seated in a wheelchair and passively transported 1.5~m forward in the dark by the experimenter. If vestibular signals alone are sufficient for path-integration, then passive transport should produce the same allocentric anchoring as active walking.

\subsection{Experiment 4: Passive Vestibular Backward Translation}

Two conditions were tested: baseline-stationary and vestibular-backward (1.5~m passive wheelchair transport backward). If the path-integration mechanism is sensitive to the direction of travel, backward transport should shift the intrinsic bias in the opposite direction compared to forward motion, causing targets to appear farther rather than nearer.

\subsection{Experiment 5: Stepladder Descent}

Nine observers stood on a stepladder (0.6~m height) in the dark and then descended to the floor before judging target locations. Two conditions were tested: stepladder (descend from 0.6~m) and ground-level baseline (stand on floor). If path-integration operates only along the horizontal direction, then vertical descent should not shift the horizontal coordinate of the intrinsic bias, but the vertical coordinate should shift downward with the observer. If path-integration is isotropic, both coordinates should remain at the stepladder position.

\section{Results}

\subsection{Experiment 1: Allocentric Coding Confirmed}

In the walking condition with a 12-second delay, observers judged target locations to be significantly nearer compared to the baseline-stationary condition (Table~\ref{tab:summary}). The average judged locations shifted approximately 1.35~m in the horizontal direction, consistent with the 1.5~m walking distance. This result is predicted by the allocentric hypothesis: the intrinsic bias remained anchored at the home base while the observer moved 1.5~m forward, so targets appeared nearer from the new vantage point.

The walking condition with a 60-second delay produced a much smaller separation from baseline compared to the 12-second delay condition, indicating that the allocentric anchoring of the intrinsic bias decays over time. With a sufficiently long delay, the intrinsic bias appears to gradually migrate toward the observer's current position.

\subsection{Experiment 2: Attention Is Required for Path-Integration}

When observers walked while performing a demanding cognitive task, the judged target locations were statistically indistinguishable from the baseline-stationary condition. The divided-attention-walking and baseline-stationary data points overlapped, and both conditions were well fitted by the same intrinsic bias curve. This result demonstrates that path-integration requires attentional resources during locomotion; when attention is diverted, the intrinsic bias travels with the observer rather than remaining anchored at the home base.

\subsection{Experiment 3: Vestibular Signals Suffice}

Passive vestibular forward translation via wheelchair produced judged target locations that were statistically indistinguishable from those obtained during active walking, and both conditions were significantly nearer than baseline-stationary. This result demonstrates that vestibular signals alone are sufficient for path-integration; active proprioceptive signals from walking are not required for the allocentric anchoring of the intrinsic bias.

\subsection{Experiment 4: Direction-Sensitive Odometer}

Passive backward wheelchair transport shifted judged target locations in the opposite direction compared to forward transport. Targets appeared farther than baseline-stationary, consistent with the intrinsic bias remaining at the home base while the observer moved backward, so that targets at a given physical distance appeared farther from the new (more proximal) vantage point. This finding confirms that the path-integration mechanism correctly encodes the direction of travel.

\subsection{Experiment 5: Horizontal-Only Path-Integration}

In the stepladder descent experiment, the horizontal coordinate of judged target locations did not differ between the stepladder and ground-level conditions, while the vertical coordinate shifted downward (Table~\ref{tab:exp5}). This pattern matches the prediction of horizontal-only path-integration: the odometer gauged only the horizontal component of the observer's displacement (which was zero, since the descent was purely vertical), leaving the horizontal coordinate of the intrinsic bias unchanged. The vertical shift is consistent with the intrinsic bias's vertical coordinate moving downward with the observer, as vertical displacement was not integrated.

\begin{table}[ht]
\centering
\caption{Summary of key findings across all five experiments.}
\label{tab:summary}
\begin{tabular}{llll}
\toprule
\textbf{Experiment} & \textbf{N} & \textbf{Manipulation} & \textbf{Key Finding} \\
\midrule
Exp 1 & 8 & Walk 1.5 m (12/60 s delay) & Allocentric coding; decays with time \\
Exp 2 & 8 & Walk with divided attention & Attention required for path-integration \\
Exp 3 & 8 & Passive forward wheelchair & Vestibular signals sufficient \\
Exp 4 & 8 & Passive backward wheelchair & Direction-sensitive odometer \\
Exp 5 & 9 & Stepladder descent (0.6 m) & Horizontal-only path-integration \\
\bottomrule
\end{tabular}
\end{table}

\begin{table}[ht]
\centering
\caption{Experiment 5 predictions and results for isotropic vs.\ horizontal-only path-integration following stepladder descent.}
\label{tab:exp5}
\begin{tabular}{lccl}
\toprule
\textbf{Hypothesis} & \textbf{Horizontal Coord.} & \textbf{Vertical Coord.} & \textbf{Observed?} \\
\midrule
Isotropic path-integration & Fixed at home base & Fixed at stepladder height & No \\
Horizontal-only path-integration & Fixed at home base & Shifts down with observer & Yes \\
\bottomrule
\end{tabular}
\end{table}

\section{Discussion}

The present findings establish that the human visual system employs an allocentric (world-centered) reference frame for coding visual space during locomotion. Across five experiments, we demonstrated that the intrinsic bias---a fundamental component of spatial vision that determines how objects are localized in the absence of visible depth cues---remains anchored at the observer's starting location rather than traveling with the observer during locomotion. This allocentric anchoring depends on a path-integration mechanism driven by vestibular signals, requires attentional resources, decays over time, is sensitive to the direction of travel, and operates anisotropically along the horizontal direction only.

\subsection{Allocentric Advantage}

The use of an allocentric reference frame for spatial perception offers a significant computational advantage during locomotion. In an egocentric system, every spatial representation would need to be continuously updated as the observer moves, requiring computationally expensive transformations for each object in the scene \citep{byrne2007}. An allocentric system, by contrast, maintains a fixed spatial reference that needs only to track the observer's displacement, reducing the computational burden to a single path-integration operation regardless of the number of objects in the environment.

\subsection{Vestibular Drive and Attention Dependency}

The finding that passive vestibular translation is sufficient for path-integration, while active proprioceptive signals are not necessary, identifies the vestibular system as the primary sensory input for the spatial odometer. This is consistent with the well-established role of vestibular signals in self-motion perception and spatial updating \citep{angelaki2008, glasauer2005}. However, the requirement for attentional resources during walking (Experiment~2) indicates that vestibular-driven path-integration is not a fully automatic process. When attentional resources are diverted to a demanding secondary task, the odometer fails to track the observer's displacement, and the intrinsic bias defaults to an egocentric mode in which it travels with the observer.

\subsection{Temporal Decay}

The observation that the allocentric anchoring effect was smaller with a 60-second delay compared to a 12-second delay suggests that the path-integration signal decays over time. This temporal dynamics may reflect the gradual accumulation of error in the vestibular-driven distance estimate, or it may indicate a time-dependent process by which the intrinsic bias is re-centered on the observer's current position when no additional motion cues are available.

\subsection{Anisotropic Path-Integration}

Perhaps the most striking finding is the horizontal-only nature of the path-integration mechanism. The stepladder descent experiment (Experiment~5) demonstrated that the odometer gauges only horizontal distance, not vertical displacement. This anisotropy has a clear parallel in the path-integration system of desert ants, which is known to operate exclusively along the horizontal plane \citep{wohlgemuth2001, wintergerst2017}. The convergence between human and ant path-integration systems suggests a shared evolutionary design principle reflecting the predominant importance of horizontal displacement during terrestrial locomotion.

\subsection{Ground-Based Spatial Coding}

Our findings are consistent with the broader framework of ground-based spatial coding, in which the ground surface serves as the primary reference for localizing objects in the visual environment \citep{ooi2001, wu2004, he2004}. The intrinsic bias itself is a manifestation of this ground-based coding system, representing an implicit assumption about the shape of the ground plane. The allocentric anchoring of this bias to a fixed ground location during locomotion reinforces the centrality of the ground surface in spatial vision.

\subsection{Limitations and Future Directions}

Several limitations should be noted. First, the experiments tested only short walking distances (1.5~m) and relatively brief delays (up to 60 seconds). The behavior of the allocentric odometer over longer distances and extended time periods remains to be characterized. Second, all experiments were conducted in darkness, which eliminates visual feedback that might modulate the path-integration mechanism. Future studies should examine how visible environmental features interact with the vestibular-driven odometer. Third, the neural substrates of the allocentric odometer remain unknown. Candidate regions include the hippocampal formation, which is known to support allocentric spatial representations \citep{okeefe1978, moser2008}, and the vestibular cortex, which processes self-motion signals.

\section{Conclusion}

We have identified a human allocentric odometer that maintains the spatial reference frame of the intrinsic bias at the observer's starting location during locomotion. This odometer is driven by vestibular signals, requires attentional resources, decays over time, is direction-sensitive, and operates exclusively along the horizontal dimension. These properties reveal a path-integration mechanism that supports stable spatial perception during self-motion and shares fundamental design principles with the horizontal odometer described in navigating insects.

\bibliographystyle{plainnat}
\bibliography{references}

\end{document}
